\documentclass[a4paper,11pt]{article}

% ---- Encoding & fonts -------------------------------------------------------
\usepackage[T1]{fontenc}
\usepackage[utf8]{inputenc}
\usepackage{lmodern}
\usepackage{microtype}              % protrusion/kerning -> fewer overfull boxes

% ---- Geometry (~1in margins) ------------------------------------------------
\usepackage[a4paper,margin=1in,headheight=15pt,headsep=14pt]{geometry}

% ---- Math -------------------------------------------------------------------
\usepackage{amsmath,amssymb}

% ---- Graphics, tables, colour ----------------------------------------------
\usepackage{graphicx}
\usepackage{booktabs}
\usepackage[table,svgnames]{xcolor}
\usepackage{adjustbox}

% ---- Callout boxes (breakable pill-tab tcolorboxes) -------------------------
\usepackage[most]{tcolorbox}
\tcbuselibrary{skins,breakable}

\usepackage{pifont}
\IfFileExists{bbding.sty}{%
  \usepackage{bbding}%
  \newcommand{\glyphHand}{\Pointinghand}%
  \newcommand{\glyphPencil}{\Pencil}%
}{%
  \newcommand{\glyphHand}{\ding{43}}%
  \newcommand{\glyphPencil}{\ding{46}}%
}

% ---- Lists & spacing --------------------------------------------------------
\usepackage{enumitem}
\setlist{leftmargin=1.5em,topsep=4pt,itemsep=3pt}

\newlist{steps}{enumerate}{1}
\setlist[steps]{label=\textbf{Step \arabic*.},
  leftmargin=4.4em, labelwidth=3.8em, labelsep=0.6em, labelindent=0pt,
  align=left, topsep=5pt, itemsep=6pt}
\usepackage{parskip}

% ---- Headers / footers ------------------------------------------------------
\usepackage{fancyhdr}

% ---- Section styling --------------------------------------------------------
\usepackage{titlesec}

% ---- Links (hidden) + URL line-breaking so long URLs never overflow ---------
\usepackage[hidelinks]{hyperref}
\usepackage{xurl}
\hypersetup{breaklinks=true}

% =============================================================================
%  COLOURS — navy chrome + the FIVE taxonomy ramps
% =============================================================================
\definecolor{navy}{HTML}{21355E}        % banner + headings
\definecolor{navyrule}{HTML}{21355E}
\definecolor{inktext}{HTML}{1A202C}     % body ink
\definecolor{mutedtext}{HTML}{6B7280}   % header / meta grey

% Intuition — blue
\definecolor{intuFrame}{HTML}{2C5AA0}   \definecolor{intuFill}{HTML}{F4F7FC}
% Everyday — amber
\definecolor{everyFrame}{HTML}{8A5A1E}  \definecolor{everyFill}{HTML}{FBF1E3}
% Worked — green
\definecolor{workFrame}{HTML}{2E7D52}   \definecolor{workFill}{HTML}{F1F8F3}
% Watch out — red
\definecolor{watchFrame}{HTML}{B23A48}  \definecolor{watchFill}{HTML}{FBF1F2}
% Key takeaway — purple/violet
\definecolor{keyFrame}{HTML}{6A4C93}    \definecolor{keyFill}{HTML}{F4F0F9}
% Wait, what? — teal
\definecolor{waitFrame}{HTML}{0F766E}   \definecolor{waitFill}{HTML}{EEF7F5}

\color{inktext}
\hypersetup{linkcolor=navy,urlcolor=navy,citecolor=navy}

\newcommand{\CourseShort}{DNN Companion}
\newcommand{\SessionNo}{12}
\newcommand{\LectureTitle}{Transformer Architectures, Positional Encoding \& BERT}

\pagestyle{fancy}
\fancyhf{}
\fancyhead[L]{\footnotesize\color{mutedtext}\textsf{\CourseShort}}
\fancyhead[R]{\footnotesize\color{mutedtext}\textsf{Session \SessionNo\ \textperiodcentered\ \LectureTitle}}
\fancyfoot[L]{\footnotesize\color{mutedtext}\textsf{WILP, BITS Pilani}}
\fancyfoot[C]{\footnotesize\color{mutedtext}\thepage}
\renewcommand{\headrulewidth}{0.4pt}
\renewcommand{\footrulewidth}{0pt}
\renewcommand{\headrule}{\hbox to\headwidth{\color{navyrule!55}\leaders\hrule height \headrulewidth\hfill}}

\titleformat{\section}
  {\normalfont\Large\bfseries\sffamily\color{navy}}
  {\thesection}{0.8em}{}
  [\vspace{2pt}\color{navyrule!40}\hrule height 0.6pt]
\titlespacing*{\section}{0pt}{18pt}{8pt}

\titleformat{\subsection}
  {\normalfont\large\bfseries\sffamily\color{navy}}
  {\thesubsection}{0.6em}{}
\titlespacing*{\subsection}{0pt}{12pt}{4pt}

\newcommand{\secsub}[1]{%
  \vspace{-6pt}{\small\sffamily\color{mutedtext}\textbf{#1}}\par\vspace{8pt}%
}

% =============================================================================
%  TCOLORBOX CALLOUT MACROS
% =============================================================================
\newtcolorbox{intuition}[1][Mathematical Intuition]{%
  enhanced, breakable, frame hidden, interior hidden,
  toptitle=5pt, bottomtitle=5pt, top=14pt, bottom=10pt, left=12pt, right=12pt,
  colback=intuFill, colframe=intuFrame,
  overlay unbroken and first={
    \draw[intuFrame, lw=1pt] (frame.south west) rectangle (frame.north east);
    \node[anchor=west, fill=intuFrame, text=white, font=\sffamily\bfseries\small,
          inner xsep=10pt, inner ysep=4pt, rounded corners=3pt]
      at ([xshift=12pt, yshift=0pt]frame.north west) {\ding{168}\ \ #1};
  }%
}

\newtcolorbox{everyday}[1][Everyday Picture]{%
  enhanced, breakable, frame hidden, interior hidden,
  toptitle=5pt, bottomtitle=5pt, top=14pt, bottom=10pt, left=12pt, right=12pt,
  colback=everyFill, colframe=everyFrame,
  overlay unbroken and first={
    \draw[everyFrame, lw=1pt] (frame.south west) rectangle (frame.north east);
    \node[anchor=west, fill=everyFrame, text=white, font=\sffamily\bfseries\small,
          inner xsep=10pt, inner ysep=4pt, rounded corners=3pt]
      at ([xshift=12pt, yshift=0pt]frame.north west) {\glyphHand\ \ #1};
  }%
}

\newtcolorbox{worked}[1][Worked Example]{%
  enhanced, breakable, frame hidden, interior hidden,
  toptitle=5pt, bottomtitle=5pt, top=14pt, bottom=10pt, left=12pt, right=12pt,
  colback=workFill, colframe=workFrame,
  overlay unbroken and first={
    \draw[workFrame, lw=1pt] (frame.south west) rectangle (frame.north east);
    \node[anchor=west, fill=workFrame, text=white, font=\sffamily\bfseries\small,
          inner xsep=10pt, inner ysep=4pt, rounded corners=3pt]
      at ([xshift=12pt, yshift=0pt]frame.north west) {\glyphPencil\ \ #1};
  }%
}

\newtcolorbox{watchout}[1][Watch Out!]{%
  enhanced, breakable, frame hidden, interior hidden,
  toptitle=5pt, bottomtitle=5pt, top=14pt, bottom=10pt, left=12pt, right=12pt,
  colback=watchFill, colframe=watchFrame,
  overlay unbroken and first={
    \draw[watchFrame, lw=1pt] (frame.south west) rectangle (frame.north east);
    \node[anchor=west, fill=watchFrame, text=white, font=\sffamily\bfseries\small,
          inner xsep=10pt, inner ysep=4pt, rounded corners=3pt]
      at ([xshift=12pt, yshift=0pt]frame.north west) {\ding{55}\ \ #1};
  }%
}

\newtcolorbox{keytake}[1][Key Takeaway]{%
  enhanced, breakable, frame hidden, interior hidden,
  toptitle=5pt, bottomtitle=5pt, top=14pt, bottom=10pt, left=12pt, right=12pt,
  colback=keyFill, colframe=keyFrame,
  overlay unbroken and first={
    \draw[keyFrame, lw=1pt] (frame.south west) rectangle (frame.north east);
    \node[anchor=west, fill=keyFrame, text=white, font=\sffamily\bfseries\small,
          inner xsep=10pt, inner ysep=4pt, rounded corners=3pt]
      at ([xshift=12pt, yshift=0pt]frame.north west) {\ding{51}\ \ #1};
  }%
}

\newtcolorbox{waitwhat}[1][Wait, What?]{%
  enhanced, breakable, frame hidden, interior hidden,
  toptitle=5pt, bottomtitle=5pt, top=14pt, bottom=10pt, left=12pt, right=12pt,
  colback=waitFill, colframe=waitFrame,
  overlay unbroken and first={
    \draw[waitFrame, lw=1pt] (frame.south west) rectangle (frame.north east);
    \node[anchor=west, fill=waitFrame, text=white, font=\sffamily\bfseries\small,
          inner xsep=10pt, inner ysep=4pt, rounded corners=3pt]
      at ([xshift=12pt, yshift=0pt]frame.north west) {\small!?\ \ #1};
  }%
}

\newcommand{\housefig}[2]{%
  \begin{figure}[htbp]
    \centering
    \includegraphics[width=0.92\linewidth]{#1}
    \caption{#2}
  \end{figure}
}

\newcommand{\flipanswer}[1]{%
  \begin{adjustbox}{angle=180}%
    \begin{minipage}{\linewidth}%
      \small\color{mutedtext}#1%
    \end{minipage}%
  \end{adjustbox}%
}

\begin{document}

% Banner Header
\begin{tcolorbox}[
  enhanced,
  colback=navy,
  colframe=navy,
  arc=0pt, outer arc=0pt,
  left=16pt, right=16pt, top=14pt, bottom=14pt
]
  \color{white}
  \begin{minipage}{0.72\linewidth}
    {\fontfamily{phv}\selectfont
      {\small\bfseries\scshape WILP, BITS Pilani \hfill Deep Neural Networks (AIML ZG511)}\par
      \vspace{4pt}
      {\LARGE\bfseries Session 12: Transformers, Positional Encoding \& BERT \par}
      \vspace{4pt}
      {\small Scaled Dot-Product Attention, Sinusoidal Embeddings, Layer Normalization, and BERT Pre-training\par}
    }
  \end{minipage}%
  \hfill
  \begin{minipage}{0.24\linewidth}
    \centering
    \color{white!80}\sffamily\scriptsize
    \textbf{Module 9}\\[2pt]
    Session 12 Reader\\[2pt]
    Full Worked Solutions
  \end{minipage}
\end{tcolorbox}

\vspace{10pt}

% Section 1
\section{The Transformer Architecture \& Self-Attention}
\secsub{Transformers eliminate recurrence completely, enabling parallel processing of entire sequences via Self-Attention.}

The Transformer model replaces sequential RNN loops with Scaled Dot-Product Self-Attention:
$$\text{Attention}(Q, K, V) = \text{softmax}\left(\frac{QK^T}{\sqrt{d_k}}\right)V$$
where $Q, K, V$ represent Query, Key, and Value projections.

\begin{worked}[Full Worked Self-Attention Numerical]
Query token 1: $q_1 = [0.21, 0.04, 0.63, 0.36]^T$. Key vectors $k_1, k_2, k_3$:
$$k_1 = \begin{bmatrix} 0.31 \\ 0.84 \\ 0.12 \\ 0.45 \end{bmatrix}, \quad k_2 = \begin{bmatrix} 0.15 \\ 0.22 \\ 0.78 \\ 0.09 \end{bmatrix}, \quad k_3 = \begin{bmatrix} 0.55 \\ 0.01 \\ 0.34 \\ 0.88 \end{bmatrix}$$
Key dimension $d_k = 4$, so scaling factor $\sqrt{d_k} = \sqrt{4} = 2.0$.

\begin{steps}
\item Compute raw dot products $q_1 \cdot k_j$:
\begin{align*}
q_1 \cdot k_1 &= 0.0651 + 0.0336 + 0.0756 + 0.1620 = 0.3363 \\
q_1 \cdot k_2 &= 0.0315 + 0.0088 + 0.4914 + 0.0324 = 0.5641 \\
q_1 \cdot k_3 &= 0.1155 + 0.0004 + 0.2142 + 0.3168 = 0.6469
\end{align*}

\item Divide by scaling factor $\sqrt{d_k} = 2.0$:
$$s_1 = 0.1682, \quad s_2 = 0.2821, \quad s_3 = 0.3235$$

\item Apply Softmax to compute attention probabilities $\alpha_{1j}$:
\begin{align*}
\exp(s_1) &= 1.1832, \quad \exp(s_2) = 1.3259, \quad \exp(s_3) = 1.3820 \\
\text{Sum} &= 1.1832 + 1.3259 + 1.3820 = 3.8911 \\
\alpha_{11} &= 0.3041, \quad \alpha_{12} = 0.3407, \quad \alpha_{13} = 0.3552
\end{align*}
\end{steps}
\end{worked}

% Section 2
\section{Positional Encoding}
\secsub{Since Self-Attention is order-agnostic, positional embeddings inject word sequence position into token vectors.}

The Transformer uses fixed sinusoidal functions for position $pos$ and dimension $i$:
$$PE_{(pos, 2i)} = \sin\left(\frac{pos}{10000^{2i/d_{\text{model}}}}\right), \quad PE_{(pos, 2i+1)} = \cos\left(\frac{pos}{10000^{2i/d_{\text{model}}}}\right)$$

\housefig{figures/positional_encoding.png}{Sinusoidal positional encoding values across token positions for different dimension frequencies.}

\begin{intuition}[Why Sinusoids?]
Sinusoidal encodings allow the model to learn relative positions easily via linear transformations.
\end{intuition}

% Section 3
\section{Normalization: LayerNorm vs BatchNorm}
\secsub{Layer Normalization normalizes across feature dimensions per sample, making it ideal for variable-length sequences.}

\housefig{figures/norm_comparison.png}{Comparison of Batch Normalization (across batch N) vs Layer Normalization (across features D).}

\begin{keytake}[LayerNorm vs BatchNorm]
\begin{center}
\begin{adjustbox}{width=\linewidth}
\begin{tabular}{lll}
\toprule
\textbf{Property} & \textbf{Batch Normalization (BatchNorm)} & \textbf{Layer Normalization (LayerNorm)} \\
\midrule
Axis & Normalizes across Batch dimension $N$ & Normalizes across Feature dimension $D$ \\
Batch Dependency & Depends heavily on batch size & Independent of batch size (works for $N=1$) \\
NLP Suitability & Poor (struggles with variable lengths) & Superior (standard in Transformers) \\
\bottomrule
\end{tabular}
\end{adjustbox}
\end{center}
\end{keytake}

% Section 4
\section{BERT \& Pre-training Paradigms}
\secsub{BERT pre-trains bidirectional encoder representations using Masked Language Modeling and Next Sentence Prediction.}

BERT (Bidirectional Encoder Representations from Transformers) uses only the Transformer Encoder stack.
Pre-training tasks include:
1. \textbf{Masked Language Model (MLM):} Mask 15\% of input tokens and predict them.
2. \textbf{Next Sentence Prediction (NSP):} Predict whether Sentence B follows Sentence A.

\begin{worked}[BERT Variants \& Specialized Models]
\begin{itemize}
  \item \textbf{BERT Base:} 12 layers, 768 hidden size, 12 heads, 110M parameters.
  \item \textbf{BERT Large:} 24 layers, 1024 hidden size, 16 heads, 340M parameters.
  \item \textbf{patentBERT / docBERT:} Fine-tuned for patent classification and document handling.
  \item \textbf{TinyBERT:} Compact student model distilled from teacher BERT for fast edge inference.
\end{itemize}
\end{worked}

% Section 5
\section{Self-Test \& Pocket Recap}
\secsub{Test your knowledge on Transformers, positional encodings, and BERT variants.}

\subsection*{Self-Test Questions}
1. \textbf{Question:} Why is scaling by $\frac{1}{\sqrt{d_k}}$ critical in dot-product attention?\\
\flipanswer{\textbf{Answer:} For large values of $d_k$, the dot product grows large in magnitude, pushing softmax into regions with extremely small gradients. Dividing by $\sqrt{d_k}$ keeps variance at $1.0$ and prevents gradient vanishing.}

\vspace{10pt}

2. \textbf{Question:} How does BERT differ from traditional decoder-only language models like GPT?\\
\flipanswer{\textbf{Answer:} BERT uses bidirectional attention (Encoder-only) to look at left and right context simultaneously, whereas GPT uses masked unidirectional attention (Decoder-only) to predict the next token.}

\end{document}
